% Generated draft: requires mathematical and source review.
\documentclass[11pt]{article}
\usepackage[margin=25mm]{geometry}
\usepackage{fontspec}
\setmainfont{Latin Modern Roman}
\usepackage{amsmath,amssymb,mathtools,unicode-math}
\setmathfont{Latin Modern Math}
\usepackage{enumitem,needspace,xurl,hyperref}
\setlength{\parindent}{0pt}
\setlength{\parskip}{5pt}
\setlength{\emergencystretch}{3em}
\newcommand{\N}{\mathbb{N}}
\newcommand{\Z}{\mathbb{Z}}
\newcommand{\Q}{\mathbb{Q}}
\newcommand{\R}{\mathbb{R}}
\newcommand{\C}{\mathbb{C}}
\newcommand{\F}{\mathbb{F}}
\newcommand{\A}{\mathbb{A}}
\newcommand{\PP}{\mathbb{P}}
\newcommand{\E}{\mathbb{E}}
\newcommand{\eps}{\varepsilon}
\newcommand{\dd}{\,\mathrm{d}}
\begin{document}
\section*{UnsolvedMath ID: 30002719}
\textbf{Source problem code:} catalogue entry
\subsection{Definitions and mathematical statement}
Under the standard univariate convention, the probability generating function of a random variable \(X\) taking values in \(\Z_{\geq 0}=\{0,1,2,\ldots\}\) is
\[
G_X(z)=\E[z^X]=\sum_{k=0}^{\infty}\Pr(X=k)z^k.
\]
This identity may be read as an identity of formal power series; the coefficients are nonnegative and sum to one.

The supplied record refers to a family \(F_n\), described as Waring-polynomial series approximating the leading root of a deformed exponential formal power series. It supplies no formula for the deformed exponential, the Waring polynomials, or \(F_n\), and does not specify the series variable, parameters, index range, or meaning of ``leading root.'' Consequently, these objects cannot be defined uniquely from the supplied evidence.

\textbf{Requested statement, conditional on the missing definitions.} Does the source-defined series \(F_n\) admit a representation
\[
F_n(z)=\E[z^{X_n}]
      =\sum_{k=0}^{\infty}\Pr(X_n=k)z^k
\]
for a nonnegative integer-valued random variable \(X_n\) with an explicit or easily describable probabilistic construction?

Here \(z\) is a placeholder for the unspecified generating-function variable, and the nonnegative integer-valued domain expresses the standard probability-generating-function convention. The supplied record does not establish that this convention matches the original source's intended variables or parameters. It also does not specify the quantification over \(n\), admissible parameter values, or a criterion for ``explicitly or easily describable.'' This is therefore an unresolved source-dependent question, not a fully specified assertion about a defined family.
\subsection{Short English statement}
The question asks whether the series \(F_n\) can be interpreted as the probability generating function of a simply described nonnegative integer-valued random variable. The supplied record does not define \(F_n\) or its parameters and index range, so a precise family-wide statement requires further source clarification.
\subsection{Sources}
\begin{itemize}
\item Supplied local UnsolvedMath registry record, permanent ID \(30002719\), title ``Probabilistic Interpretation of Waring-Polynomial Series,'' source problem code \texttt{catalogue entry}, in the Oberwolfach Reports Open Problems collection. The record gives proposed year 2014 and identifies no individual proposer.
\item Supplied source citation: \emph{Probability, Trees and Algorithms (2015)}, Oberwolfach Reports, pp. 2825--2871, DOI \texttt{10.4171/owr/2014/50}. \url{https://doi.org/10.4171/owr/2014/50}. The supplied background describes Oberwolfach Reports as documenting mathematical workshops and their open-problem discussions. The citation's parenthetical year 2015 and the record's proposed year 2014 are retained without reconciliation.
\item Supplied literature source: ``Kuznetsov, deformed exponential zeros.'' \url{https://arxiv.org/abs/2412.02462}. No more detailed bibliographic information or supporting passage is supplied.
\item Supplied literature source: ``Oberwolfach report.'' \url{https://doi.org/10.4171/owr/2014/50}.
\item Supplied literature-check date: \texttt{2026-08-21}. The supplied dated curation labels the problem open and states that only small cases and computational evidence are known for the probability-generating-function question, with no general probabilistic representation for \(F_n\) established. This is the supplied assessment, not an independently verified literature conclusion. The background explicitly calls for independent verification of time-sensitive status and cited claims.
\item The registry labels the statement \texttt{corrected\_verified} and reports manual semantic, typographic, and LaTeX review asserting self-containment and meaning preservation. These are supplied review claims; they do not supply the missing definitions or establish mathematical correctness.
\end{itemize}
\end{document}
